\section{Singularity extraction on the touching cells}
\label{sec:extraction}

\propTitle{Non-smooth terms of the Green's function}{taylor} For $R\ge0$,
\begin{equation}
    \label{eq:taylor}
    g(R) = \frac{\mathrm{e}^{-\I kR}}{R}
         = \sum_{n=0}^{\infty} \frac{(-\I k)^n}{n!}\,R^{n-1}
         = \sum_{\ell=0}^{\infty} \frac{(-1)^\ell k^{2\ell}}{(2\ell)!}\,R^{2\ell-1}
         \;-\;\I \sum_{\ell=0}^{\infty} \frac{(-1)^\ell k^{2\ell+1}}{(2\ell+1)!}\,R^{2\ell}.
\end{equation}
The series converge absolutely and locally uniformly for $R>0$; the second one is an entire function of $R^2$. As
functions of $\V{z}=\V{r}-\V{r}'\in\Rn^2$ the even powers $R^{2\ell}=\NORM{\V{z}}^{2\ell}$ are polynomials, while the odd powers
$R^{2\ell-1}$ are not smooth at $\V{z}=\V{0}$.

\vspace{\baselineskip}
\proof Insert the exponential series and separate the even and odd $n$; this is Proposition~2.2
in Section~2.6 of~\cite{machacek2025method}. \qed

\defTitle{Extraction of order $N$}{extraction} For $N\in\N$ let
\begin{equation}
    \label{eq:sn}
    S_N(R) = \sum_{\ell=0}^{N-1} \frac{(-1)^\ell k^{2\ell}}{(2\ell)!}\,R^{2\ell-1}
           = \frac{1}{R} - \frac{k^2}{2}R + \frac{k^4}{24}R^3 - \dots, \qquad
    g_N = g - S_N .
\end{equation}
Table~\ref{tab:extraction} lists the orders: $N=1$ is the classical extraction of the static singularity~\cite{Tzoulis2004}, $N=2$ the two-term extraction
of~\cite{machacek2025method} (Entry~2.1 in Section~2.6), there written as
\begin{equation*}
    g=g_{\mathrm{s}}+\frac{1}{4\uppi}R^{-1}-\frac{k^2}{8\uppi}R,
\end{equation*}
and $N=3$ the one of this note.
For every $N\ge1$, $g_N$ extends continuously to $R=0$ with $g_N(0)=-\I k$.

\lemmaTitle{Regularity of the remainder}{regularity} For $N\ge1$ the function $\V{z}\mapsto g_N(\NORM{\V{z}})$ is
real-analytic on $\Rn^2\setminus\{\V{0}\}$ and belongs to $C^{2N-2,1}$ on every bounded set: its derivatives up to the
order $2N-2$ are continuous, and those of the order $2N-2$ are Lipschitz.
It is not of class $C^{2N-1}$ unless $k=0$.

\vspace{\baselineskip}
\proof By~\eqref{eq:taylor},
\begin{equation*}
    g_N(R) = E(R^2) + R^{2N-1}F(R^2),
\end{equation*}
with the entire functions
\begin{equation*}
    E(s) = -\I\sum_{\ell\ge0}\frac{(-1)^\ell k^{2\ell+1}}{(2\ell+1)!}s^\ell, \qquad
    F(s) = \sum_{\ell\ge N}\frac{(-1)^\ell k^{2\ell}}{(2\ell)!}s^{\ell-N}.
\end{equation*}
The function $E(\NORM{\V{z}}^2)$ is real-analytic on $\Rn^2$.

The function $\NORM{\V{z}}^{2N-1}=\bigl(\NORM{\V{z}}^2\bigr)^{N-1}\NORM{\V{z}}$ is smooth away from the origin and
positively homogeneous of degree $2N-1$. Hence its derivatives of order $m\le 2N-2$ are homogeneous of degree
$2N-1-m\ge1$. They are therefore continuous at the origin (with the value zero), and those of order $2N-2$, being
homogeneous of degree one, are Lipschitz. The derivatives of order $2N-1$ are homogeneous of degree zero and not constant,
so they are bounded but discontinuous at the origin.

Multiplication by the analytic factor $F(\NORM{\V{z}}^2)$ preserves the class $C^{2N-2,1}$ on bounded sets. Since
$F(0)=(-1)^Nk^{2N}/(2N)!\ne0$ for $k\ne0$, the term $F(0)\NORM{\V{z}}^{2N-1}$ is present and the class is sharp. \qed

\vspace{\baselineskip}
On a touching cell pair the integrand of a moment meets $\V{z}=\V{0}$, in the interior of the four-dimensional domain for
the coincident cell and on its boundary for a shared edge or vertex. A tensor \GL{} rule integrates polynomials exactly,
so its error on such a pair is governed by the regularity of Lemma~\lemmaRef{regularity} (Table~\ref{tab:extraction}).

\begin{table}[ht]
    \centering
    \caption{The extraction orders and the regularity of the remainder on a touching cell pair.}
    \label{tab:extraction}
    \renewcommand{\arraystretch}{1.25}
    \begin{tabular}{clll}
        \toprule
        $N$ & extracted $S_N$ & used in & remainder $g_N$ \\ \midrule
        $0$ & --- & & unbounded \\
        $1$ & $1/R$ & \cite{Tzoulis2004} & Lipschitz \\
        $2$ & $1/R-k^2R/2$ & \cite{machacek2025method} & $C^{2,1}$ \\
        $3$ & $1/R-k^2R/2+k^4R^3/24$ & this note & $C^{4,1}$ \\ \bottomrule
    \end{tabular}
\end{table}

The error then decreases algebraically in the order $q$, and the faster, the larger $N$. The lemma is not a quadrature
error estimate; for that theory see~\cite{sauter2011boundary}. On the cell pairs that do not touch, $g$ is
real-analytic on a neighbourhood of the domain, and the \GL{} error decreases exponentially.

\coroTitle{Exact decomposition of the touching moments}{decomposition} For every $N\ge1$, $\offs\in\mathcal{T}$ and
$\alpha,\beta\in\{0,1\}$,
\begin{equation}
    \label{eq:decomposition}
    M^{[g]}_{\alpha\beta}(\offs) = M^{[g_N]}_{\alpha\beta}(\offs)
        + \sum_{\ell=0}^{N-1}\frac{(-1)^\ell k^{2\ell}}{(2\ell)!}\,\Mom{2\ell-1}{\alpha\beta}{\offs},
\end{equation}
and likewise for the $N$-moments. The first term is evaluated with \GL{}; the static moments $\Mom{p}{\alpha\beta}{\offs}$,
$p\in\{-1,1,3\}$, are given in closed form by Theorem~\theoRef{touching}.

\vspace{\baselineskip}
\proof The moments are linear in $K$, and every term is absolutely integrable: the kernels $R^{2\ell-1}$ with $\ell\ge0$ have
$2\ell-1\ge-1>-2$, and $g_N$ is bounded. \qed

\vspace{\baselineskip}
\noindent\textbf{What has to be computed.} By Lemma~\lemmaRef{symmetries}, the static moments of the four offsets
\begin{equation*}
    \offs\in\{(0,0),\,(1,0),\,(0,1),\,(1,1)\}
\end{equation*}
determine all nine touching offsets. These are the coincident cell, the two kinds of shared edge and the shared vertex.
For $N=3$ we need, for each of them, the four $x$-moments $M_{00},M_{10},M_{01},M_{11}$ and the three $y$-moments,
with $p\in\{-1,1,3\}$. The $y$-moments follow from the $x$-moments by Lemma~\lemmaRef{symmetries}(iii).
